\section{Theoretical analysis}\label{sec:theory}

The individual tools used below---split-conformal calibration, the union bound, order statistics---are standard, and so are most of the arguments. What this section contributes is an account of what their composition guarantees to an agent that must decide under irreversibility with a stochastic oracle, and of the conditions under which it does. Propositions~\ref{prop:sound} and~\ref{prop:residual} make the certifier/refuter/residual decomposition exact. Theorem~\ref{thm:bound} applies conformal calibration with a Bonferroni correction to the statistic the fence actually thresholds---a maximum over a variable-size candidate set behind decidable vetoes---and makes its validity condition explicit: the pool must be drawn from the rule's own jurisdiction, and a pool that is not can hide its own violation (Section~\ref{sec:res-bound}). Theorem~\ref{thm:tight} records that the classical bound is attained by that statistic, so the correction and the budget are forced rather than conservative. Proposition~\ref{prop:dilemma} and Corollary~\ref{cor:mondrian} state the refinement dilemma exactly. The remaining propositions delimit what is guaranteed; they are stated because nearby designs violate them.

\subsection{Assumptions and their status}\label{sec:assumptions}

Table~\ref{tab:assumptions} lists every assumption and how the experiments satisfy or test it. Several are environmental, and a deployment that does not supply them does not obtain the corresponding guarantee.

{\footnotesize
\begin{longtable}{@{}p{3.1cm}p{2.4cm}p{7.6cm}@{}}
\caption{Assumptions, their type, and their status in the implementation and experiments. ``Measured'' entries test an assumption on this corpus; they do not establish it in general.}\label{tab:assumptions}\\
\toprule
Assumption & Type & Status in this work \\
\midrule
\endfirsthead
\toprule
Assumption & Type & Status in this work \\
\midrule
\endhead
A1 Stable run identifier across retries and replays & environment & Supplied by the harness; a retry under a new identifier is out of scope. \\
A2 Deterministic canonical identity: $\kappa_t$ and $H$ are pure; $H$ is collision-free & implementation & Pure canonicalisation; $H$ is SHA-256; unit-tested. \\
A3 Atomic insert-if-absent on $\mathit{op}$ & storage & \emph{Partial:} unique for a single writer but not crash-atomic; crashes injected at protocol points, not inside a write. \\
A4 Effect before commit; commit only on success & implementation & Algorithm~\ref{alg:fence}; eight injected crash points (Section~\ref{sec:res-crash}). \\
A5 Every declared field policy is semantics-preserving & schema & Tool author's obligation. Measured: 0 of 232 distinct test pairs canonically equal. \\
A5$'$ A difference in a declared typed field implies a different act (one direction only) & schema & Tool author's obligation. Measured: 0 of 281 duplicates refuted; 20 of 20 typed-field injections refuted. \\
A5$''$ Projection completeness: $\sigma_t(x)=\sigma_t(x')\Leftrightarrow x\approx_t x'$ & schema (optional) & Only in Proposition~\ref{prop:compose}, Case A; false for every tool in the corpus. \\
A6 Conditionally on $m$, each score of a truly distinct residual comparison is exchangeable with the calibration pool & statistical & Stratified pair-level split, attacks held out; 82.8\% of test pairs share a seed act with a calibration pair, tested by seed-level resplitting (Section~\ref{sec:res-bound}). Needs a stationary judge; unverifiable at deployment. \\
A7 No concurrent dispatch per identity & scheduler & Sequential dispatch in all experiments; concurrency not tested. \\
A8 The provider applies at most one effect per idempotency key & provider & Only in Proposition~\ref{prop:compose}; absent by design in the crash experiment. \\
\bottomrule
\end{longtable}}

Structurally, the candidate set of a request is every effect committed in the same run for the same tool; duplicates across runs or across tools are outside every result below.

\subsection{Soundness of the decidable layers}

\begin{proposition}[Decidable layers]\label{prop:sound}
Under A2 and A5, every certifier suppression is correct: if the fence suppresses on the canonical branch, the request denotes the same act as the matched committed effect. Under A5$'$, no judge output can cause suppression of a request whose declared typed fields differ from those of every committed effect.
\end{proposition}

\begin{proof}
The canonical branch fires only if $H(\kappa_t(x))=H(\kappa_t(x'))$ for a committed $x'$; by A2 this implies $\kappa_t(x)=\kappa_t(x')$ and by A5 $x\approx_t x'$. For the second claim, a committed effect whose projection differs is removed from $C$ before any judge call, so it is never the maximising candidate; if every committed effect differs, $C=\emptyset$ and the fence executes.
\end{proof}

\subsection{A bound on the deployed statistic}\label{sec:bound}

The fence suppresses on $S^\ast=\max_{c\in C}J(t,x,c)$, not on one pairwise score, and a pairwise bound $P(s_i>\hat t)\le\alpha$ does not imply $P(\max_i s_i>\hat t)\le\alpha$. The result below is stated for the deployed statistic.

\begin{theorem}[Multiplicity-corrected wrong-suppression bound]\label{thm:bound}
Let the pool $\{s_i\}_{i\le n}$ be drawn from the rule's jurisdiction (Definition~\ref{def:jurisdiction}) and let $\thi(m)$ be given by Eq.~\eqref{eq:thi}. Suppose (A6) that, conditionally on $m$, the judge score of any comparison between a request and a candidate that denotes a different act is exchangeable with the pool. Then for every request $x$ that denotes a different act from every member of its candidate set $C$, $|C|=m$,
\[
P\big(S^\ast(x,C)>\thi(m)\big)\le\alpha ,
\]
for every $m\ge0$, with no assumption on the distribution of scores and no assumption on the dependence among the $m$ comparisons.
\end{theorem}

\begin{proof}
Condition on $m$. If $m=0$ or $\thi(m)=+\infty$ the fence cannot suppress and the probability is $0$. Otherwise let $\beta=\alpha/m$ and $k=\lceil(n+1)(1-\beta)\rceil\le n$, so $\thi(m)=s_{(k)}$. Fix a candidate $c_i\in C$. Because $x\not\approx_t c_i$, its score $s'_i$ is a distinct-pair score, so by A6 the $n+1$ values $s_1,\dots,s_n,s'_i$ are exchangeable and, with ties broken uniformly at random, the rank of $s'_i$ among them is uniform on $\{1,\dots,n+1\}$. If $s'_i>s_{(k)}$, then $s'_i$ exceeds at least $k$ pool values and has rank at least $k+1$ under any tie-breaking, so
\[
P\big(s'_i>s_{(k)}\big)\le\frac{n+1-k}{n+1}\le 1-(1-\beta)=\beta .
\]
This holds for each $i$ separately. The union bound over the $m$ candidates, which requires no independence, gives
\[
P\big(S^\ast>\thi(m)\big)=P\big(\exists i\le m: s'_i>\thi(m)\big)\le\sum_{i=1}^m P\big(s'_i>\thi(m)\big)\le m\beta=\alpha .
\]
The bound holds conditionally on each value of $m$ and therefore also after averaging over the distribution of $m$. Abstentions only remove candidates from the maximum while $\thi(m)$ is computed with the full $m$, which is conservative.
\end{proof}

The probability is joint over the calibration pool and the request, so the bound is \emph{marginal}: a particular deployment's pool can exceed $\alpha$ (Section~\ref{sec:res-bound}), and nothing is claimed under distribution shift, including injection (Section~\ref{sec:limitations}). It tolerates arbitrary dependence among the $m$ comparisons, which share $x$.

\paragraph{The pool must match the jurisdiction}
A6 compares the pool with the scores the deployed rule thresholds, and which pairs those are is decided by the layers in front of the judge. A pool drawn from a wider population satisfies A6 only by coincidence; if the extra pairs score lower, the threshold falls and the bound fails while the rate over a whole test split still looks acceptable. Jurisdiction matching is therefore a necessary condition of the theorem, not an implementation detail.

\paragraph{The escalation floor}
A duplicate's candidate set contains its match, so taking a maximum helps the escalation side and $\tlo$ needs no correction. With integer-percent scores ties are frequent and $\tlo$ is capped at $\thi$, so we treat the band as an empirical quantity (Section~\ref{sec:res-components}) and claim no bound for it.

\subsection{Calibration budget and the refinement dilemma}\label{sec:dilemma}

\begin{proposition}[Calibration budget; refinement dilemma]\label{prop:dilemma}
Let a configuration have pool size $n$ and level $\alpha$.
\begin{enumerate}[(a),leftmargin=*,itemsep=1pt]
\item $\thi(m)$ is finite if and only if $m\le\alpha(n+1)$; equivalently, semantic suppression at candidate-set size $m$ requires $n\ge\lceil m/\alpha\rceil-1$, and the largest certifiable candidate-set size is $m^\ast=\lfloor\alpha(n+1)\rfloor$.
\item When finite, $\thi(m)$ is the $\rho$-th \emph{largest} pool score with $\rho=\lfloor (n+1)\alpha/m\rfloor$, which is non-decreasing in $n$.
\item If $(\kappa',\sigma')$ refines $(\kappa,\sigma)$ and both are calibrated on the same calibration split, the refined pool is a sub-multiset of the original pool. Hence $n'\le n$, $m^{\ast\prime}\le m^\ast$ and $\rho'\le\rho$: refinement can only shrink the certifiable budget and move the threshold to an equal or higher rank from the top.
\end{enumerate}
\end{proposition}

\begin{proof}
(a) Since $n$ is an integer, $\lceil(n+1)(1-\alpha/m)\rceil\le n$ holds iff $(n+1)(1-\alpha/m)\le n$, i.e.\ iff $m\le\alpha(n+1)$. (b) The $k$-th smallest of $n$ values is the $(n+1-k)$-th largest, and $n+1-\lceil(n+1)(1-\alpha/m)\rceil=\lfloor (n+1)\alpha/m\rfloor$. (c) The pool is the set of distinct calibration pairs in the jurisdiction with a score; the jurisdiction of the proposed configuration is the residual, and by Proposition~\ref{prop:residual}(iii) the refined residual is contained in the original one.
\end{proof}

At $\alpha=0.05$ and $m=1$, the declared pool of the primary judge ($n=142$) places $\thi$ at the 7th largest score, the tightened pool of Section~\ref{sec:res-refinement} ($n=49$) at the 2nd largest, and any pool with $n\le18$ at $+\infty$. The proposition does \emph{not} say the refined threshold is numerically higher; that depends on whether the judge scores the removed pairs lower than those that remain, which Section~\ref{sec:res-refinement} measures. Asking for the bound to hold within each kind of difference is a second refinement---of the calibration population rather than of the schema---and it pays the same price.

\begin{corollary}[Class-conditional budget]\label{cor:mondrian}
Partition the jurisdiction into classes $c=1,\dots,K$ by a function of the request, let the pool contain $n_c$ scores of class $c$ ($\sum_c n_c=n$), and let $t_c(m)$ be Eq.~\eqref{eq:thi} computed on the class-$c$ pool alone.
\begin{enumerate}[(a),leftmargin=*,itemsep=1pt]
\item If A6 holds within each class, then $P(S^\ast>t_c(m)\mid\text{class } c)\le\alpha$ for every $c$.
\item $t_c(m)$ is finite iff $m\le\alpha(n_c+1)$. Every class is certifiable at $m$ iff $\min_c n_c\ge\lceil m/\alpha\rceil-1$, which requires $n\ge K(\lceil m/\alpha\rceil-1)$; and $\min_c\lfloor\alpha(n_c+1)\rfloor\le\lfloor\alpha(n/K+1)\rfloor$.
\item If the class is not observable when the fence decides, thresholding at $\max_c t_c(m)$ satisfies (a) for every class simultaneously; it is finite only if every class is certifiable at $m$.
\end{enumerate}
\end{corollary}

Each class pool is a sub-multiset of the pool, so its threshold sits at an equal or higher rank from the top: partitioning the jurisdiction into $K$ classes gives each class the budget of Proposition~\ref{prop:dilemma}(a) at its own $n_c$. A per-class guarantee costs calibration budget, and when the classes are defined by the label, as the kinds of difference in our corpus are, it costs the deduplication of the most conservative class (Section~\ref{sec:res-class}).

\subsection{The standard bound is attained}\label{sec:tight}

Theorem~\ref{thm:bound} rests on a union bound. That the union bound is sharp when nothing is known about the dependence among the events is classical: $P(\bigcup_iA_i)\le\min\{1,\sum_iP(A_i)\}$ is the upper Fr\'echet bound, attained when the events are disjoint or cover the space~\cite{frechet1935,ruschendorf1982maxsums,galambos1996bonferroni}, and the sharper bounds of Hunter and Worsley need the pairwise intersection probabilities~\cite{hunter1976union,worsley1982bonferroni}, which a fence does not know. In the language of merging p-values, the Bonferroni rule is precise and admissible under arbitrary dependence~\cite{vovk2020averaging,vovk2022admissible}. We therefore do not present the tightness below as a new technical result. It answers a narrower question: is the classical worst case reachable by the \emph{deployed} statistic under the fence's own assumption, and does any rule a fence can audit do better? Within deterministic rules that do not depend on the scale of the judge's scores---which are uncalibrated, so a rule that reads meaning into their absolute values inherits that miscalibration---the standard bound is exactly attained. Lemma~\ref{lem:scalefree}, which identifies the class, is a characterisation of a known type---symmetric functions that commute with every increasing bijection are order statistics~\cite{ovchinnikov1997invariant,marichal2002order}---stated here without a continuity assumption and proved in \ref{app:proofs} for completeness.

\begin{lemma}[Scale-free thresholds]\label{lem:scalefree}
Let $\tau$ map pools of $n$ distinct scores to $[-\infty,+\infty]$, be symmetric in its arguments, and satisfy $\tau(g(s_1),\dots,g(s_n))=g(\tau(s_1,\dots,s_n))$ for every strictly increasing bijection $g$ of $\mathbb{R}$, extended by $g(\pm\infty)=\pm\infty$. Then $\tau\equiv s_{(k)}$ for one fixed $k\in\{1,\dots,n\}$, or $\tau\equiv+\infty$, or $\tau\equiv-\infty$.
\end{lemma}

A deterministic rule that thresholds $S^\ast$ at a scale-free function of the pool is therefore ``suppress iff $S^\ast>s_{(k)}$'' for a fixed $k$, or never suppresses, or always does. The fence is such a rule, with $k=k^\ast=\lceil(n+1)(1-\alpha/m)\rceil$ when $k^\ast\le n$ and $\rho=n+1-k^\ast$ as in Proposition~\ref{prop:dilemma}. Let $\mathcal{Q}_{n,m}$ be the class of joint laws of $(s_1,\dots,s_n,s'_1,\dots,s'_m)$ in which, for each $i$, the pool and $s'_i$ are $n+1$ independent draws from one continuous distribution, with arbitrary dependence among $s'_1,\dots,s'_m$. Every law in $\mathcal{Q}_{n,m}$ satisfies A6, so a rule that is valid under A6 must be valid on all of $\mathcal{Q}_{n,m}$.

\begin{theorem}[The standard bound is attained by the deployed statistic]\label{thm:tight}
For every $k\in\{1,\dots,n\}$,
\[
\sup_{\mathcal{Q}_{n,m}} P\Big(\max_{i\le m}s'_i>s_{(k)}\Big)=\mathbb{E}\big[\min\{1,mB\}\big],\qquad B\sim\mathrm{Beta}(n+1-k,\,k),
\]
and the supremum is attained. For the fence's threshold,
\[
\sup_{\mathcal{Q}_{n,m}}P(\text{wrong suppression})\;=\;\frac{m\rho}{n+1}-\mathbb{E}\big[(mB-1)_+\big]\;\le\;\frac{m\rho}{n+1}\;\le\;\alpha,
\]
where $\mathbb{E}[(mB-1)_+]\le(m-1)\exp\!\big(-(1-\alpha)^2n/(2m)\big)$.
\end{theorem}

The attaining law is not exotic: pool and candidate scores are i.i.d., the most favourable reading of A6, and only the joint behaviour of one request's $m$ candidate scores is adversarial---spread so that at most one is high, as when a request is compared with committed effects that differ from it in different ways.

\begin{corollary}[What cannot be improved]\label{cor:tight}
Fix $\alpha\in(0,1)$ and $m\ge1$. Worst cases are over $\mathcal{Q}_{n,m}$.
\begin{enumerate}[(a),leftmargin=*,itemsep=1pt]
\item \emph{The constant.} For $m=1$ the worst case is $\rho/(n+1)$, which equals $\alpha$ whenever $(n+1)\alpha$ is an integer; for every $m$ it tends to $\alpha$ as $n\to\infty$. No constant below $\alpha$ bounds wrong suppression for all pool sizes.
\item \emph{The correction.} For $m=1$, every rank $k<k^\ast$ has worst case $(n+1-k)/(n+1)>\alpha$. For every $m$, a rule with a fixed per-candidate level $\gamma>\alpha/m$ has worst case tending to $\min\{1,m\gamma\}>\alpha$. The Bonferroni level is the largest per-candidate level valid for every pool size without an assumption on the dependence among the $m$ comparisons.
\item \emph{The budget.} If $m>\alpha(n+1)$, the most conservative finite threshold, the pool maximum ($k=n$), has worst case exactly
\[
\mathbb{E}\big[\min\{1,mB\}\big]=\frac{m}{n+1}\Big(1-\Big(1-\frac1m\Big)^{n+1}\Big),\qquad B\sim\mathrm{Beta}(1,n),
\]
and every other finite threshold has a larger one. This exceeds $\alpha$ unless $m/(n+1)$ lies within $\delta_{0.05}<1.1\times10^{-10}$ of $\alpha$ (at $\alpha=0.05$; proof in \ref{app:proofs}), which no pool with $n+1<4.8\times10^8$ can do. There, switching semantic suppression off beyond the budget is forced rather than conservative: the only valid deterministic scale-free rule is never to suppress.
\end{enumerate}
\end{corollary}

Proofs are in \ref{app:proofs}. For every configuration in this paper the gap between worst case and bound is below $10^{-10}$: 4.895\% for the declared pool ($n=142$, $m=1$), 4.000\% for the tightened pool ($n=49$), and 4.975\% at $n=200$ for $m=1,2,5,10$, the simulation setting of Section~\ref{sec:res-bound}, whose law gave 3.7\% at $m=10$. One rank lower the worst case exceeds $\alpha$ (5.594\% and 6.000\%). Under the attaining law the uncorrected rule suppresses 49.7\% of distinct requests at $m=10$ and 99.9\% at $m=50$, against 38.3\% and 89.0\% in the simulation. A direct simulation of the attaining law agrees with the exact values in all eight configurations we computed; the budget condition was checked for every $n+1\le10^6$ at $\alpha=0.05$ and $\alpha=0.01$ with no exception; and the equality in Corollary~\ref{cor:tight}(c) was checked without its closed form, against the general rank-$k$ computation at $k=n$ to within $10^{-12}$ over 4,000 pairs $(n,m)$, and by simulation, which agrees in all six budget configurations (supplementary material, Section~S21). The result concerns deterministic rules: a coin-tossing rule beyond the budget is valid but cannot be audited or replayed. With integer-percent scores, ties lower the realised rate (Section~\ref{sec:res-components}).

\subsection{Why the architecture is layered}

\begin{proposition}[No zero-error threshold on overlapping scores]\label{prop:impossible}
Let $F_{\mathrm{same}}$ and $F_{\mathrm{diff}}$ be the judge-score distributions of duplicate and distinct pairs. If they overlap---some score has positive probability under both---no rule that thresholds the score alone achieves zero wrong suppression and zero silent escape simultaneously. An escalation band can remove silent escape only by flagging, not suppressing, the duplicates in the overlap.
\end{proposition}

The proof is in \ref{app:proofs}. Zero error on a class needs a decision that does not consult the score, which is what the certifier and refuter are; the overlap is judge- and tool-dependent (supplementary material, Section~S8).

\subsection{Fence-level at-most-once}\label{sec:amo}

\begin{proposition}[Fence-level at-most-once commitment]\label{prop:amo}
Under A1--A4 and A7, for a fixed run $r$ and operation identifier $\mathit{op}$, across any number of serialised executions and process restarts sharing the durable log, the effect for $\mathit{op}$ is dispatched and committed at most once by the fence, except that a crash strictly between the external effect and the commit insert leaves a \textsc{pending} intent with no commit row, after which a later execution may dispatch again.
\end{proposition}

The proof is in \ref{app:proofs}. The commit window cannot be closed by any local mechanism, because during it the only record of the effect is at the provider. Without A7, two concurrent executions can both miss the log and dispatch; at most one commit row survives and the loser records \textsc{commit\_conflict}, so the duplicate is observable but not prevented.

\subsection{External exactly-once only by composition}\label{sec:compose}

\begin{proposition}[Composition with a provider key]\label{prop:compose}
\emph{Case A.} If a key function $\pi_t$ is paraphrase-invariant ($x\approx_t x'\Rightarrow\pi_t(x)=\pi_t(x')$) and complete for identity ($\pi_t(x)=\pi_t(x')\Rightarrow x\approx_t x'$), and the provider satisfies A8, then presenting $\pi_t(x)$ as the provider's idempotency key yields external exactly-once application for that tool, including across the commit window. A schema satisfying A5$''$ supplies $\pi_t=\sigma_t$.
\emph{Case B.} If the schema satisfies A5$'$ but not A5$''$, then $\sigma_t$ is not admissible as a provider key (distinct acts sharing a projection would be merged at the provider, a wrong suppression that no local audit can see), and the canonical identifier $\mathit{op}$ is admissible but not paraphrase-invariant; the commit window then remains open for paraphrased retries.
\end{proposition}

\begin{proposition}[Identity recognition is not substitutable]\label{prop:structure}
Mis-recognised identity defeats both fence-level commitment and a correct provider key: an unrecognised duplicate presents a fresh key, which a correct provider must apply, and a wrongly recognised distinct act presents a stale key, which a correct provider absorbs. Conversely, a provider key satisfying Case A of Proposition~\ref{prop:compose} makes A7 unnecessary for at-most-once application.
\end{proposition}

Proofs are in \ref{app:proofs}. Identity recognition (A) is thus the one guarantee neither the log nor the provider can supply, because both are keyed on its output.

\subsection{What is proved and what is measured}

Proved: exactness of the decidable layers under A2, A5 and A5$'$ (Proposition~\ref{prop:sound}); the marginal bound under A6 with its budget, overall and per class (Theorem~\ref{thm:bound}, Proposition~\ref{prop:dilemma}, Corollary~\ref{cor:mondrian}); attainment of the standard bound by the deployed statistic (Theorem~\ref{thm:tight}, Corollary~\ref{cor:tight}); fence-level at-most-once under A1--A4 and A7 (Proposition~\ref{prop:amo}); external exactly-once only by composition (Proposition~\ref{prop:compose}). Everything else is measured: whether the assumptions hold on this corpus, the residual's size, each judge's deduplication, how the realised rate varies across pools, splits and kinds of difference, behaviour under injection, and the composition part of the refinement dilemma.
